% Dyadic chain-level arithmetization: square classes, syndromes, decoding, moments, LDPC.

\paragraph{胞腔链的平方类算术实现与边界综合征}\label{par:spg-dyadic-squareclass-syndrome}
在固定 dyadic 立方复形中，边界链不仅可由素因子集合编码为整数（定理 \ref{thm:spg-boundary-godelization-holographic-dictionary}），且其模 \(2\) 的链复形结构可完全由平方类群的乘法结构实现。
该口径把 \(\partial\circ\partial=0\) 译为“综合征为 \(1\)”的局部闭合约束，并给出显式的线性时间 bulk 译码与由单一边界整数读出任意单项式体积分的表列方案。

\begin{theorem}[素数平方类实现的胞腔链复形：边界--算术同构]\label{thm:spg-dyadic-squareclass-chain-complex}
固定 \(m\ge 1\) 并记 \(\mathcal Q_m\) 为 \(I^n=[0,1]^n\) 的 \(m\)-dyadic 立方剖分所诱导的 \(n\)-维立方复形。
对每一维 \(k\in\{0,\dots,n\}\) 取定取向的 \(k\)-胞腔集合 \(\mathcal F_k(\mathcal Q_m)\)。
对每个 \(k\) 选取一族两两不同的素数 \(\{p_\sigma:\sigma\in\mathcal F_k(\mathcal Q_m)\}\)，并令
\[
\mathcal G_k:=\Big\langle [p_\sigma]\ :\ \sigma\in\mathcal F_k(\mathcal Q_m)\Big\rangle\ \le\ \QQ^\times/(\QQ^\times)^2
\]
为其生成的平方类群，其中 \([r]\) 表示 \(r\in\QQ^\times\) 在 \(\QQ^\times/(\QQ^\times)^2\) 中的平方类。
定义映射
\[
\mathbf G_k:\ C_k(\mathcal Q_m;\FF_2)\longrightarrow \mathcal G_k,\qquad
\mathbf G_k\Bigl(\sum_{\sigma\in\mathcal F_k(\mathcal Q_m)}\varepsilon_\sigma\,\sigma\Bigr)
:=\prod_{\sigma\in\mathcal F_k(\mathcal Q_m)} [p_\sigma]^{\varepsilon_\sigma},
\]
其中 \(\varepsilon_\sigma\in\{0,1\}\)。
则 \(\mathbf G_k\) 为 \(\FF_2\)-线性同构；并且存在唯一群同态 \(\delta_k:\mathcal G_k\to\mathcal G_{k-1}\) 使得对任意 \(\sigma\in\mathcal F_k(\mathcal Q_m)\) 有
\[
\delta_k([p_\sigma])=\prod_{\tau\prec\sigma}[p_\tau]^{[\tau:\sigma]\bmod 2},
\]
其中 \([\tau:\sigma]\in\{0,\pm 1\}\) 为胞腔边界的符号系数，\(\tau\prec\sigma\) 表示 \((k-1)\)-胞腔 \(\tau\) 为 \(\sigma\) 的面。
此外交换关系成立：
\[
\mathbf G_{k-1}\circ \partial_k=\delta_k\circ \mathbf G_k.
\]
\end{theorem}
\begin{proof}
\(\mathbf G_k\) 将链群 \(C_k(\mathcal Q_m;\FF_2)\cong \FF_2^{|\mathcal F_k|}\) 的标准基元 \(\sigma\) 送至平方类基元 \([p_\sigma]\)，故为同构。
对基元 \(\sigma\)，有
\(
\partial_k\sigma=\sum_{\tau\prec\sigma}[\tau:\sigma]\tau
\)，模 \(2\) 后系数为 \([\tau:\sigma]\bmod 2\)，从而
\[
\mathbf G_{k-1}(\partial_k\sigma)
=\prod_{\tau\prec\sigma}[p_\tau]^{[\tau:\sigma]\bmod 2}
=\delta_k([p_\sigma]).
\]
对任意链作线性延拓即得交换性。
唯一性由 \(\mathcal G_k\) 由 \([p_\sigma]\) 生成且乘法关系仅来自指数模 \(2\) 加法给出。证毕。
\end{proof}

\begin{proposition}[平方类支撑距离与 Hamming 距离的一致性]\label{prop:spg-squareclass-support-hamming-isometry}
沿用定理 \ref{thm:spg-dyadic-squareclass-chain-complex} 的记号，固定 \(k\in\{0,\dots,n\}\)。
对 \(H\in\mathcal G_k\) 定义其素支撑
\[
\mathrm{supp}(H):=\{\sigma\in\mathcal F_k(\mathcal Q_m):\ [p_\sigma]\ \text{在 }H\text{ 的平方自由代表中出现}\}.
\]
并定义平方类支撑距离
\[
d_{\square}(H_1,H_2):=\card{\mathrm{supp}(H_1H_2^{-1})}.
\]
则对任意 \(c_1,c_2\in C_k(\mathcal Q_m;\FF_2)\) 有
\[
\boxed{\;
d_{\square}\bigl(\mathbf G_k(c_1),\mathbf G_k(c_2)\bigr)
\ =\
d_H(c_1,c_2)\;},
\]
其中 \(d_H\) 为系数向量的 Hamming 距离。
\end{proposition}
\begin{proof}
写 \(c_i=\sum_{\sigma}\varepsilon^{(i)}_\sigma\,\sigma\)（\(\varepsilon^{(i)}_\sigma\in\{0,1\}\)）。
由定义
\[
\mathbf G_k(c_1)\mathbf G_k(c_2)^{-1}
=\prod_{\sigma}[p_\sigma]^{\varepsilon^{(1)}_\sigma-\varepsilon^{(2)}_\sigma}
=\prod_{\sigma}[p_\sigma]^{\varepsilon^{(1)}_\sigma\oplus\varepsilon^{(2)}_\sigma}.
\]
因此 \(\mathrm{supp}(\mathbf G_k(c_1)\mathbf G_k(c_2)^{-1})=\{\sigma:\varepsilon^{(1)}_\sigma\neq \varepsilon^{(2)}_\sigma\}\)，其基数即为 Hamming 距离 \(d_H(c_1,c_2)\)。证毕。
\leanverified{paper\_spg\_squareclass\_support\_hamming\_isometry}
\end{proof}

\begin{corollary}[边界 Gödel 比值的平方类等距性]\label{cor:spg-boundary-godel-ratio-squareclass-isometry}
取 \(k=n-1\)。
对 dyadic polycube \(U,V\subset I^n\)，令 \(H_m(U),H_m(V)\) 为定理 \ref{thm:spg-boundary-godelization-holographic-dictionary} 的边界 Gödel 整数，并以 \([\cdot]\) 表示其在 \(\QQ^\times/(\QQ^\times)^2\) 中的平方类。
定义平方类比值
\[
R_m(U,V):=[H_m(U)]\,[H_m(V)]^{-1}.
\]
则有
\[
\boxed{\;
d_{\square}\bigl([H_m(U)],[H_m(V)]\bigr)
\ =\
\card{\mathrm{supp}(R_m(U,V))}
\ =\
d_H\bigl(\partial_n[U],\partial_n[V]\bigr)\; }.
\]
\leanverified{paper\_spg\_boundary\_godel\_ratio\_squareclass\_isometry}
\end{corollary}
\begin{proof}
定理 \ref{thm:spg-boundary-godelization-holographic-dictionary} 的编码恰对应于定理 \ref{thm:spg-dyadic-squareclass-chain-complex} 在 \(k=n-1\) 的情形：\([H_m(U)]=\mathbf G_{n-1}(\partial_n[U])\)。
因此由命题 \ref{prop:spg-squareclass-support-hamming-isometry} 立得。证毕。
\end{proof}

\begin{theorem}[边界 Gödel 码的算术化差分距离与 Stokes 读出的 Lipschitz 稳定性]\label{thm:spg-godel-boundary-gcd-lipschitz-stability}
沿用定理 \ref{thm:spg-boundary-godelization-holographic-dictionary} 的边界 Gödel 编码并固定 \(m,n\ge 1\)。
记 \(\omega(N)\) 为正整数 \(N\) 的互异素因子个数。
对任意 dyadic polycube \(U,V\subset I^n\)，定义 Gödel 对称差计数
\[
\operatorname{Dist}_m(U,V)
:=\omega\!\left(\frac{H_m(U)\,H_m(V)}{\gcd\!\bigl(H_m(U),H_m(V)\bigr)^2}\right).
\]
则成立：
\begin{enumerate}
  \item \textbf{（边界差的面积等式）}
        若将 \(\partial U-\partial V\) 视为取向 \((n-1)\)-面上的整数链并令
        \[
        \operatorname{Area}(\partial U-\partial V)
        :=2^{-m(n-1)}\bigl\|\partial_n[U]-\partial_n[V]\bigr\|_{\ell^1},
        \]
        则有精确恒等式
        \[
        \operatorname{Area}(\partial U-\partial V)=2^{-m(n-1)}\,\operatorname{Dist}_m(U,V).
        \]
  \item \textbf{（体积读出的 Lipschitz 控制）}
        令 \(\omega_0\) 为 \eqref{eq:spg-universal-stokes-probe-form} 中满足 \(\dd\omega_0=\dd x_1\cdots\dd x_n\) 的通用 Stokes 探针，则
        \[
        \bigl|\operatorname{Vol}(U)-\operatorname{Vol}(V)\bigr|
        =\left|\int_{\partial U-\partial V}\omega_0\right|
        \le 2^{-m(n-1)}\,\operatorname{Dist}_m(U,V).
        \]
\end{enumerate}
\end{theorem}
\begin{proof}
由于 \(H_m(U)\) 与 \(H_m(V)\) 皆为平方自由整数，其最大公因子 \(\gcd(H_m(U),H_m(V))\) 的素因子集合恰为两者边界支撑的交；
因而
\(
H_m(U)H_m(V)/\gcd(H_m(U),H_m(V))^2
\)
的素因子集合恰为两者边界支撑的对称差。
另一方面，\(\partial_n[U]-\partial_n[V]\) 在每个取向 \((n-1)\)-面上的系数属于 \(\{-1,0,1\}\)，其 \(\ell^1\) 范数正是该对称差的基数。
由推论 \ref{cor:spg-boundary-godel-ratio-squareclass-isometry}，该基数也等于 \(\card{\mathrm{supp}(R_m(U,V))}\)，从而
\(
\|\partial_n[U]-\partial_n[V]\|_{\ell^1}=\operatorname{Dist}_m(U,V)
\)。
乘以单个 dyadic \((n-1)\)-面的面积 \(2^{-m(n-1)}\) 即得第一式。
\par
第二式由 Stokes 读出恒等式 \(\operatorname{Vol}(W)=\int_{\partial W}\omega_0\)（见 \eqref{eq:spg-probe-volume-identity}）给出
\[
\operatorname{Vol}(U)-\operatorname{Vol}(V)=\int_{\partial U-\partial V}\omega_0.
\]
又由 \eqref{eq:spg-universal-stokes-probe-form} 可取 \(\|\omega_0\|_\infty\le 1\)，故对任意取向面 \(\vec f\) 有
\(
\bigl|\int_{\vec f}\omega_0\bigr|\le \int_{\vec f}1\,\dd S=2^{-m(n-1)}
\)。
将 \(\partial U-\partial V\) 分解为取向面之和并用三角不等式即得
\[
\left|\int_{\partial U-\partial V}\omega_0\right|
\le 2^{-m(n-1)}\bigl\|\partial_n[U]-\partial_n[V]\bigr\|_{\ell^1}
=2^{-m(n-1)}\,\operatorname{Dist}_m(U,V).
\]
证毕。
\end{proof}

\begin{corollary}[Gödel--Stokes 边界距离下的通量维读出稳定性]\label{cor:spg-godel-stokes-distance-dimension-stability}
令 \(A\subset I^n\) 为紧集，\(U_m(A)\) 为定理 \ref{thm:spg-dyadic-outer-approx-stokes-gain-minkowski-readout} 的 dyadic 外近似。
令 \((V_m)_{m\ge 1}\) 为任意 dyadic polycube 序列，并沿用定理 \ref{thm:spg-godel-boundary-gcd-lipschitz-stability} 的 \(\omega_0\) 与 \(\operatorname{Dist}_m\) 记号。
设存在 \(d\in(0,n]\) 与常数 \(C_->0\) 使对充分大 \(m\) 有
\[
\operatorname{Vol}\bigl(U_m(A)\bigr)\ \ge\ C_-\,2^{-m(n-d)}.
\]
若进一步满足
\[
\operatorname{Dist}_m\bigl(U_m(A),V_m\bigr)=o\!\bigl(2^{m(d-1)}\bigr),
\]
则体积/通量读出在该误差下渐近稳定：
\[
\lim_{m\to\infty}\left|\frac{\int_{\partial V_m}\omega_0}{\int_{\partial U_m(A)}\omega_0}-1\right|=0,
\]
并且
\[
\limsup_{m\to\infty}\frac{\log\!\left(\int_{\partial V_m}\omega_0\right)}{m\log 2}
=
\limsup_{m\to\infty}\frac{\log\!\left(\int_{\partial U_m(A)}\omega_0\right)}{m\log 2}.
\]
因此，单一探针的 dyadic 维读出
\(
d_A:=n+\limsup_{m\to\infty}\frac{\log\!\left(\int_{\partial U_m(A)}\omega_0\right)}{m\log 2}
\)
在 \(V_m\) 上保持不变：
\[
d_A
=
n+\limsup_{m\to\infty}\frac{\log\!\left(\int_{\partial V_m}\omega_0\right)}{m\log 2}.
\]
\end{corollary}
\begin{proof}
由定理 \ref{thm:spg-godel-boundary-gcd-lipschitz-stability}(2) 及 \eqref{eq:spg-probe-volume-identity}，
\[
\left|\int_{\partial V_m}\omega_0-\int_{\partial U_m(A)}\omega_0\right|
=\bigl|\operatorname{Vol}(V_m)-\operatorname{Vol}(U_m(A))\bigr|
\le 2^{-m(n-1)}\operatorname{Dist}_m\bigl(U_m(A),V_m\bigr).
\]
两边除以 \(\operatorname{Vol}(U_m(A))\ge C_-2^{-m(n-d)}\) 得
\[
\left|\frac{\operatorname{Vol}(V_m)}{\operatorname{Vol}(U_m(A))}-1\right|
\le \frac{1}{C_-}\,\operatorname{Dist}_m\bigl(U_m(A),V_m\bigr)\,2^{-m(d-1)}
\xrightarrow[m\to\infty]{}0,
\]
从而得到通量比值收敛到 \(1\)。
由于 \(\log(1+o(1))=o(1)\)，对数化并除以 \(m\log 2\) 即得两条 \(\limsup\) 恒等式。证毕。
\end{proof}

\begin{corollary}[边界 Gödel 码的局部综合征刻画：可填充当且仅当局部闭合，且填充唯一]\label{cor:spg-dyadic-boundary-syndrome-fillability}
沿用定理 \ref{thm:spg-dyadic-squareclass-chain-complex} 的记号，并取 \(k=n-1\)。
令 \(H\in\mathcal G_{n-1}\) 对应 \((n-1)\)-链 \(b=\mathbf G_{n-1}^{-1}(H)\in C_{n-1}(\mathcal Q_m;\FF_2)\)。
则以下条件等价：
\begin{enumerate}
  \item 存在 \(x\in C_n(\mathcal Q_m;\FF_2)\) 使得 \(b=\partial_n x\)；
  \item \(\delta_{n-1}(H)=1\)，等价地 \(\partial_{n-1}b=0\)。
\end{enumerate}
并且在上述等价条件成立时，满足 \(b=\partial_n x\) 的 \(x\) 唯一。
\end{corollary}
\begin{proof}
由定理 \ref{thm:spg-dyadic-squareclass-chain-complex} 的交换性，
\[
\delta_{n-1}(H)=1
\iff
\mathbf G_{n-2}(\partial_{n-1}b)=1
\iff
\partial_{n-1}b=0,
\]
其中最后一步用 \(\mathbf G_{n-2}\) 为同构。
若 \(b=\partial_n x\)，则 \(\partial_{n-1}b=\partial_{n-1}\partial_n x=0\)，故 (1)\(\Rightarrow\)(2)。
反之，\(|\mathcal Q_m|=I^n\) 可缩，故 \(H_{n-1}(\mathcal Q_m;\FF_2)=0\) 且任意 \((n-1)\)-循环皆为 \(\partial_n\) 的像，故 (2)\(\Rightarrow\)(1)。
唯一性由 \(\partial_n\) 的单射性（定理 \ref{thm:spg-dyadic-cubical-boundary-injective}）推出。证毕。
\end{proof}

\begin{theorem}[线性时间 bulk 译码：由边界素因子分解显式重构 dyadic polycube]\label{thm:spg-linear-time-bulk-decode-from-boundary}
设 \(U\subset I^n\) 为 dyadic polycube，并以 \(\FF_2\) 系数识别其占据链 \([U]\in C_n(\mathcal Q_m;\FF_2)\)。
令 \(\Gamma\) 为 \(\mathcal Q_m\) 的 \(n\)-胞腔邻接图：顶点为全部 \(n\)-胞腔并附加一顶点 \(\infty\) 表示外部；若两 \(n\)-胞腔共用一张 \((n-1)\)-面，则连一条边；对贴边 \(n\)-胞腔，再以其外侧 \((n-1)\)-面与 \(\infty\) 连边。
由边界链 \(\partial_n[U]\) 在 \(\Gamma\) 上定义边标记 \(\eta:E(\Gamma)\to\FF_2\)：若对应的 \((n-1)\)-面在 \(\partial_n[U]\) 的支撑中出现，则 \(\eta(e)=1\)，否则 \(\eta(e)=0\)。
则存在显式算法在时间 \(O(|V(\Gamma)|+|E(\Gamma)|)\) 内重构 \(U\)：
固定 \(x(\infty)=0\)，任取生成树 \(T\subset\Gamma\) 并沿树边递推赋值
\[
x(v)=x(u)+\eta(uv)\quad(\bmod 2),
\]
从而得到唯一 \(x\in C_n(\mathcal Q_m;\FF_2)\) 满足 \(\partial_n x=\partial_n[U]\)，且该 \(x\) 等于 \([U]\)。
\end{theorem}
\begin{proof}
对任意 \(x\in C_n(\mathcal Q_m;\FF_2)\)，其边界满足：若 \(uv\) 为内部相邻边，则 \(\eta(uv)=x(u)+x(v)\)；若为外边 \(v\infty\)，则 \(\eta(v\infty)=x(v)\)。
因此在生成树上由 \(x(\infty)=0\) 递推得到的 \(x\) 在每条树边满足对应约束。
需证该赋值与路径选择无关。
由 \(\partial_{n-1}\partial_n[U]=0\) 得到 \(\partial_n[U]\) 为 \((n-1)\)-循环；在图论语言下等价于任一回路上 \(\sum_{e\in\mathrm{cycle}}\eta(e)=0\)（\(\FF_2\)）。
于是沿回路递推无矛盾，故所得 \(x\) 与生成树选择无关并满足全体边约束，进而满足 \(\partial_n x=\partial_n[U]\)。
最后由 \(\partial_n\) 的单射性（定理 \ref{thm:spg-dyadic-cubical-boundary-injective}）可知满足该边界的顶维链唯一，故 \(x=[U]\)。证毕。
\end{proof}

\begin{theorem}[整系数体--界回放的线性时间算法与不一致证书]\label{thm:spg-integer-bulk-decode-or-cycle-certificate}
令 \(\partial_n:C_n(\mathcal Q_m;\ZZ)\to C_{n-1}(\mathcal Q_m;\ZZ)\) 为 dyadic \(n\)-立方复形上的边界算子，并令 \(\Gamma\) 为定理 \ref{thm:spg-linear-time-bulk-decode-from-boundary} 中的 \(n\)-胞腔邻接图（含外部顶点 \(\infty\)）。
则存在确定性算法：输入任意 \((n-1)\)-链 \(b\in C_{n-1}(\mathcal Q_m;\ZZ)\)，在时间 \(O(|E(\Gamma)|)\) 内输出下列二者之一：
\begin{enumerate}
  \item 唯一的 \(n\)-链 \(c\in C_n(\mathcal Q_m;\ZZ)\) 使 \(\partial_n c=b\)；
  \item 一条显式闭回路 \(\gamma\subset\Gamma\)（由 \((n-1)\)-面对应的边组成），满足
  \[
  \sum_{e\in\gamma}\operatorname{sgn}_\gamma(e)\,\eta_b(e)\neq 0,
  \]
  从而推出 \(b\notin \mathrm{im}(\partial_n)\)。
\end{enumerate}
其中 \(\eta_b:E(\Gamma)\to\ZZ\) 为由 \(b\) 诱导的边标记（见证明中的定义），\(\operatorname{sgn}_\gamma(e)\in\{\pm1\}\) 为回路行进方向与边定向的一致性符号。
此外，在情形 (1) 中所得 \(c\) 唯一性由 \(\partial_n\) 的单射性（定理 \ref{thm:spg-dyadic-cubical-boundary-injective}）保证。
\end{theorem}
\begin{proof}
对每条边 \(e\in E(\Gamma)\)（对应某个有向 \((n-1)\)-面 \(\vec f\)），令 \(\eta_b(e)\in\ZZ\) 为 \(b\) 在该有向面上的系数；对外边 \(v\infty\) 亦同样以外侧面取向定义 \(\eta_b(v\infty)\)。
设 \(c=\sum_{Q}u(Q)\,Q\in C_n(\mathcal Q_m;\ZZ)\)（\(Q\) 遍历 \(n\)-胞腔）为未知体链。
由立方复形的定向约定，若 \(e=(Q\to Q')\) 为内部相邻边（两胞腔共用同一 \((n-1)\)-面），则对应面在 \(\partial_n Q\) 与 \(\partial_n Q'\) 中符号相反，从而
\[
(\partial_n c)(e)=u(Q)-u(Q').
\]
若 \(e=(Q\to\infty)\) 为外边，则
\(
(\partial_n c)(e)=u(Q)
\)。
因此方程 \(\partial_n c=b\) 等价于在图 \(\Gamma\) 上求势函数 \(u:V(\Gamma)\to\ZZ\)（约定 \(u(\infty)=0\)）使得对每条有向边 \(e=(x\to y)\) 满足
\[
u(x)-u(y)=\eta_b(e).
\]
\par
取 \(\Gamma\) 的一棵生成树 \(T\)，以 \(\infty\) 为根；沿 \(T\) 自根递推定义 \(u\)：若 \(e=(x\to y)\in T\) 且 \(u(x)\) 已定，则令 \(u(y):=u(x)-\eta_b(e)\)。
该过程访问每条树边一次，时间 \(O(|T|)\)。
\par
随后对每条非树边 \(e=(x\to y)\notin T\) 检验一致性。
若发现 \(u(x)-u(y)\neq \eta_b(e)\)，则令 \(\gamma\) 为由树中唯一路径 \(x\rightsquigarrow y\) 与边 \(e\) 拼成的基本回路，并按回路行进方向赋予 \(\operatorname{sgn}_\gamma(\cdot)\)。
沿回路求和得到
\[
\sum_{e'\in\gamma}\operatorname{sgn}_\gamma(e')\,\eta_b(e')
=u(x)-u(y)-\eta_b(e)\neq 0,
\]
从而 \(\gamma\) 为不一致证书；由于任意势函数的环路差分和必须为 \(0\)，可知无解，亦即 \(b\notin\mathrm{im}(\partial_n)\)。
\par
若所有非树边均通过检验，则差分约束在全图成立，故由 \(u\) 构造的 \(c=\sum_Q u(Q)Q\) 满足 \(\partial_n c=b\)。
总时间为树递推 \(O(|T|)\) 加非树边检验 \(O(|E(\Gamma)|)\)，即 \(O(|E(\Gamma)|)\)。
最后，若存在解则由定理 \ref{thm:spg-dyadic-cubical-boundary-injective} 的单射性知解唯一。证毕。
\end{proof}

\begin{theorem}[边界 Gödel 整数的多项式矩全读出：边界一整数确定任意单项式体积分]\label{thm:spg-boundary-godel-moment-readout}
设 \(U\subset I^n\) 为 dyadic polycube。
对任意多重指标 \(\alpha=(\alpha_1,\dots,\alpha_n)\in\NN^n\)，定义 \((n-1)\)-形式
\[
\omega_\alpha
:=\frac1n\sum_{i=1}^n (-1)^{i-1}\frac{x^{\alpha+e_i}}{\alpha_i+1}\,
\dd x_1\wedge\cdots\wedge\widehat{\dd x_i}\wedge\cdots\wedge \dd x_n,
\]
其中 \(e_i\) 为第 \(i\) 个标准基。
则有恒等式
\[
\int_U x^\alpha\,\dd x_1\cdots \dd x_n
=\int_{\partial U}\omega_\alpha.
\]
进一步，若以两两不同的素数 \(\{q_{\vec f}\}\) 对 \(\partial U\) 的有向 \((n-1)\)-面 \(\vec f\) 作平方自由编码
\(
H_m(U)=\prod_{\vec f} q_{\vec f}^{\varepsilon_{\vec f}}
\)
（定理 \ref{thm:spg-boundary-godelization-holographic-dictionary}），则存在仅依赖于 dyadic 坐标与取向的权重
\[
w_\alpha(\vec f):=\int_{\vec f}\omega_\alpha
\]
使得
\[
\int_{\partial U}\omega_\alpha=\sum_{\vec f:\ \varepsilon_{\vec f}=1} w_\alpha(\vec f).
\]
因此，\(H_m(U)\) 的素因子支撑确定任意单项式矩 \(\int_U x^\alpha\).
\end{theorem}
\begin{proof}
直接计算可得 \( \dd\omega_\alpha=x^\alpha\,\dd x_1\cdots \dd x_n \)。
由 Stokes 定理 \(\int_U \dd\omega_\alpha=\int_{\partial U}\omega_\alpha\) 即得第一式。
第二式由积分的线性性与 \(\partial U\) 的面分解推出；而 \(\varepsilon_{\vec f}\) 的支撑与 \(H_m(U)\) 的素因子支撑一一对应。证毕。
\end{proof}

\begin{theorem}[固定分辨率 dyadic polycube 的有限矩完备性]\label{thm:spg-dyadic-finite-moment-completeness}
固定 \(m,n\ge 1\)，并令 \(N:=2^m\)。
设 \(U\subset I^n=[0,1]^n\) 为分辨率 \(m\) 的 dyadic polycube，即
\[
U=\bigsqcup_{\mathbf{k}\in S} C_{\mathbf{k}},
\qquad
C_{\mathbf{k}}:=\prod_{i=1}^{n}\Bigl[\frac{k_i}{N},\frac{k_i+1}{N}\Bigr],
\qquad
\mathbf{k}\in\{0,\dots,N-1\}^{n},
\]
其中 \(S\subseteq\{0,\dots,N-1\}^{n}\)。
对每个多重指标 \(\alpha=(\alpha_1,\dots,\alpha_n)\in\{0,\dots,N-1\}^{n}\)，定义单项式矩
\[
\mu_{\alpha}(U):=\int_{U}x^{\alpha}\,\dd x
=\int_U x_1^{\alpha_1}\cdots x_n^{\alpha_n}\,\dd x.
\]
则映射
\[
U\longmapsto \bigl(\mu_{\alpha}(U)\bigr)_{\alpha\in\{0,\dots,N-1\}^{n}}
\]
在全体分辨率 \(m\) 的 dyadic polycube 上为单射。
换言之，固定分辨率下的有限矩族
\(
\{\mu_{\alpha}(U):0\le \alpha_i\le N-1\}
\)
已经唯一决定 \(U\)。
\end{theorem}
\begin{proof}
令 \(\mathcal V_m\) 为所有在网格 \(\{C_{\mathbf{k}}\}\) 上分段常数的可积函数空间，则
\[
\mathcal V_m
=
\operatorname{span}\{\ind_{C_{\mathbf{k}}}:\mathbf{k}\in\{0,\dots,N-1\}^{n}\},
\qquad
\dim \mathcal V_m=N^n.
\]
定义线性映射
\[
\mathcal M_m:\mathcal V_m\to \RR^{\{0,\dots,N-1\}^{n}},
\qquad
\mathcal M_m(f):=\left(\int_{I^n}f(x)x^\alpha\,\dd x\right)_{\alpha\in\{0,\dots,N-1\}^{n}}.
\]
只需证明 \(\mathcal M_m\) 为单射。

对基元 \(\ind_{C_{\mathbf{k}}}\) 而言，
\[
\int_{I^n}\ind_{C_{\mathbf{k}}}(x)x^\alpha\,\dd x
=\prod_{i=1}^{n}\int_{k_i/N}^{(k_i+1)/N}x_i^{\alpha_i}\,\dd x_i.
\]
因此，在张量积基下，\(\mathcal M_m\) 的矩阵为 \(M^{\otimes n}\)，其中一维矩阵 \(M\in M_N(\RR)\) 由
\[
M_{j,k}:=\int_{k/N}^{(k+1)/N}x^j\,\dd x,
\qquad
0\le j,k\le N-1
\]
给出。
于是只需证明 \(M\) 可逆。

设 \(c=(c_0,\dots,c_{N-1})^{\mathsf T}\in\ker M\)。
则对一切 \(0\le j\le N-1\) 有
\[
0=\sum_{k=0}^{N-1}c_k\int_{k/N}^{(k+1)/N}x^j\,\dd x.
\]
两边同乘 \((j+1)N^{j+1}\)，得到
\[
0=\sum_{k=0}^{N-1}c_k\bigl((k+1)^{j+1}-k^{j+1}\bigr),
\qquad
0\le j\le N-1.
\]
定义多项式空间上的线性泛函
\[
\mathcal L_c(p):=\sum_{k=0}^{N-1}c_k\bigl(p(k+1)-p(k)\bigr).
\]
由上式可知 \(\mathcal L_c(t^\ell)=0\) 对 \(\ell=1,\dots,N\) 成立，故 \(\mathcal L_c\) 在所有满足 \(\deg p\le N\) 且 \(p(0)=0\) 的多项式上恒为零。

另一方面，对此类多项式直接整理得
\[
\mathcal L_c(p)
=
c_{N-1}p(N)+\sum_{k=1}^{N-1}(c_{k-1}-c_k)p(k).
\]
考虑评价映射
\[
E:\{p\in\RR[t]:\deg p\le N,\ p(0)=0\}\to \RR^N,
\qquad
E(p):=(p(1),\dots,p(N)).
\]
若 \(E(p)=0\)，则 \(p\) 在 \(0,1,\dots,N\) 共 \(N+1\) 个点为零，而 \(\deg p\le N\)，故 \(p\equiv 0\)。
因此 \(E\) 为同构。
由于 \(\mathcal L_c(p)\) 对所有此类 \(p\) 恒为零，作为 \((p(1),\dots,p(N))\) 的线性泛函，其系数必须全为零，即
\[
c_{N-1}=0,
\qquad
c_{k-1}-c_k=0\quad(1\le k\le N-1).
\]
从而 \(c_0=\cdots=c_{N-1}=0\)，故 \(\ker M=\{0\}\)，即 \(M\) 可逆。

于是 \(M^{\otimes n}\) 可逆，从而 \(\mathcal M_m\) 为线性同构。
特别地，\(\mathcal M_m\) 在 dyadic polycube 的示性函数子类上为单射，结论成立。
\end{proof}

\begin{corollary}[边界 Gödel 整数的有限矩全息完备性]\label{cor:spg-boundary-godel-finite-moment-completeness}
固定 \(m,n\ge 1\)，并令 \(N:=2^m\)。
对任意分辨率 \(m\) 的 dyadic polycube \(U\subset I^n\)，其边界 Gödel 整数 \(H_m(U)\) 在固定分辨率类内唯一决定 \(U\)。
更精确地，由定理 \ref{thm:spg-boundary-godel-moment-readout}，\(H_m(U)\) 决定全部单项式矩 \(\mu_{\alpha}(U)\)；而由定理 \ref{thm:spg-dyadic-finite-moment-completeness}，仅需有限矩盒
\[
\bigl\{\mu_{\alpha}(U):\ \alpha\in\{0,\dots,N-1\}^{n}\bigr\}
\]
即已足以唯一恢复 \(U\)。
\end{corollary}
\begin{proof}
定理 \ref{thm:spg-boundary-godel-moment-readout} 表明，边界 Gödel 整数 \(H_m(U)\) 的素因子支撑确定任意单项式体积分 \(\mu_{\alpha}(U)\)。
因此它尤其确定有限矩盒
\(
\{\mu_{\alpha}(U):\alpha\in\{0,\dots,N-1\}^{n}\}
\)。
再由定理 \ref{thm:spg-dyadic-finite-moment-completeness} 的单射性，得到 \(H_m(U)\) 在分辨率 \(m\) 的 dyadic polycube 类内唯一决定 \(U\)。
\end{proof}

\begin{remark}[与低阶通量矩障碍的相容性]\label{rem:spg-dyadic-finite-moment-completeness-vs-obstruction}
定理 \ref{thm:spg-dyadic-finite-moment-completeness} 与定理 \ref{thm:spg-prouhet-thue-morse-obstruction-dyadic-polyclube-flux-moments} 并不矛盾。
后者讨论的是总次数上界 \(|\alpha|\le r\) 固定时的低阶通量矩族，其信息量不随分辨率增长；
而前者采用的是随分辨率 \(m\) 扩张的张量积矩盒 \(0\le \alpha_i\le 2^m-1\)，其坐标数恰与分段常数空间 \(\mathcal V_m\) 的维数同阶，因而在固定分辨率类上达到线性完备。
\end{remark}

\paragraph{线性矩全息的最小维数与总次数门槛}
在固定分辨率 dyadic polycube 的张量积矩盒完备性与 Prouhet--Thue--Morse 型低阶障碍已经确立之后，还可把线性矩全息的必要维数条件进一步压缩为下列两条结论。

\begin{theorem}[固定分辨率 dyadic polycube 的线性矩全息最小维数]\label{thm:spg-linear-moment-holography-minimal-dimension}
固定 \(m,n\ge 1\)，并令 \(N:=2^m\)。沿用定理 \ref{thm:spg-dyadic-finite-moment-completeness} 的记号，记
\[
\mathcal V_m
=
\operatorname{span}\{\ind_{C_{\mathbf{k}}}:\mathbf{k}\in\{0,\dots,N-1\}^{n}\},
\qquad
\dim \mathcal V_m=N^n=2^{mn}.
\]
对每个多重指标 \(\alpha\in\NN^n\)，定义线性矩泛函
\[
\mu_\alpha(f):=\int_{I^n}f(x)x^\alpha\,\dd x
\qquad (f\in \mathcal V_m).
\]
再定义张量积矩盒
\[
\mathcal M_m^{\Box}
:=
\{\mu_\alpha:\ 0\le \alpha_i\le N-1\ \text{对全部 }i\},
\]
以及总次数截断矩族
\[
\mathcal M_r^{\triangle}
:=
\{\mu_\alpha:\ |\alpha|\le r\}.
\]
则成立：
\begin{enumerate}
  \item 对任意有限线性泛函族
  \[
  \Lambda_m=\{\ell_1,\dots,\ell_L\}\subset \mathcal V_m^\ast,
  \]
  若测量映射
  \[
  \mathcal E_{\Lambda_m}:\mathcal V_m\longrightarrow \RR^L,
  \qquad
  f\longmapsto (\ell_1(f),\dots,\ell_L(f))
  \]
  在整个 \(\mathcal V_m\) 上单射，则必有
  \[
  L\ge 2^{mn}.
  \]
  \item 张量积矩盒 \(\mathcal M_m^{\Box}\) 的坐标数恰为
  \[
  |\mathcal M_m^{\Box}|=N^n=2^{mn},
  \]
  并且其对应测量映射在线性空间 \(\mathcal V_m\) 上为同构；因此 \(\mathcal M_m^{\Box}\) 达到上述下界，是固定分辨率 dyadic polycube 在线性矩全息意义下的最小维数实现。
  \item 当 \(r=m-1\) 时，总次数截断矩族 \(\mathcal M_r^{\triangle}\) 甚至在 dyadic polycube 的示性函数子类上都不是单射，因而更不可能在整个 \(\mathcal V_m\) 上单射。
\end{enumerate}
\leanverified{paper\_spg\_linear\_moment\_holography\_minimal\_dim\_seeds}
\end{theorem}
\begin{proof}
对第 \(1\) 条，由线性代数的秩--零化度公式，
\[
\dim\ker \mathcal E_{\Lambda_m}
=
\dim \mathcal V_m-\rank(\mathcal E_{\Lambda_m})
\ge
\dim \mathcal V_m-L.
\]
若 \(L<\dim \mathcal V_m=2^{mn}\)，则 \(\ker \mathcal E_{\Lambda_m}\neq 0\)，故 \(\mathcal E_{\Lambda_m}\) 不可能单射。于是任何线性单射测量都必须满足
\[
L\ge 2^{mn}.
\]

对第 \(2\) 条，坐标数计算为
\[
|\mathcal M_m^{\Box}|=N^n=2^{mn}.
\]
而定理 \ref{thm:spg-dyadic-finite-moment-completeness} 的证明已经表明，由这些矩所组成的映射
\[
\mathcal M_m:\mathcal V_m\to \RR^{\{0,\dots,N-1\}^{n}}
\]
是线性同构。故 \(\mathcal M_m^{\Box}\) 以恰好 \(2^{mn}\) 个标量观测达到第 \(1\) 条中的必要下界，从而确为最小维数实现。

对第 \(3\) 条，把定理 \ref{thm:spg-prouhet-thue-morse-obstruction-dyadic-polyclube-flux-moments} 应用于 \(r=m-1\)。该定理给出两个互异的 \(m\)-级 dyadic polycube \(U,V\subset I^n\)，使得对一切 \(|\alpha|\le m-1\) 都有
\[
M_\alpha(U)=M_\alpha(V).
\]
而在该定理的证明中已经使用散度定理给出
\[
M_\alpha(\Omega)=(n+|\alpha|)\int_\Omega x^\alpha\,\dd x
=(n+|\alpha|)\mu_\alpha(\ind_\Omega).
\]
由于 \(n+|\alpha|>0\)，遂得
\[
\mu_\alpha(\ind_U)=\mu_\alpha(\ind_V)
\qquad (|\alpha|\le m-1).
\]
特别地，这说明总次数截断矩族在示性函数
\[
\ind_U,\ \ind_V\in \mathcal V_m
\]
上已无法分离，故在整个 \(\mathcal V_m\) 上更不可能单射。证毕。
\end{proof}

\begin{corollary}[各向同性总次数矩的线性完备门槛具有 \texorpdfstring{$2^m$}{2^m} 级尺度]\label{cor:spg-total-degree-moment-threshold-exponential-scale}
在定理 \ref{thm:spg-linear-moment-holography-minimal-dimension} 的设定下，若总次数截断矩族
\[
\mathcal M_r^{\triangle}
=
\{\mu_\alpha:\ |\alpha|\le r\}
\]
要满足“在线性维数上可能完备”的必要条件，则必须有
\[
|\mathcal M_r^{\triangle}|=\binom{n+r}{n}\ge 2^{mn}.
\]
因此当 \(n\) 固定而 \(m\to\infty\) 时，任意可能达到线性完备门槛的次数参数都必须满足
\[
r\ge (n!)^{1/n}\,2^m\,(1+o(1)).
\]
换言之，各向同性总次数门槛的真实尺度是 \(2^m\) 级，而不是 \(m\) 级。
\leanverified{paper\_spg\_total\_degree\_moment\_threshold}
\end{corollary}
\begin{proof}
总次数不超过 \(r\) 的 \(n\) 元单项式数目为
\[
|\mathcal M_r^{\triangle}|=\binom{n+r}{n}.
\]
若该矩族要在整个 \(\mathcal V_m\) 上具备线性单射的可能性，则由定理 \ref{thm:spg-linear-moment-holography-minimal-dimension} 第 \(1\) 条，必要条件必为
\[
\binom{n+r}{n}\ge 2^{mn}.
\]

现令 \(n\) 固定并令 \(m\to\infty\)。由于右端趋于无穷，满足该必要条件的 \(r=r(m)\) 也必趋于无穷。对固定 \(n\)，标准渐近给出
\[
\binom{n+r}{n}
=
\frac{r^n}{n!}\left(1+O_n(r^{-1})\right)
\qquad (r\to\infty).
\]
因此必要条件推出
\[
\frac{r^n}{n!}\left(1+O_n(r^{-1})\right)\ge 2^{mn}.
\]
取 \(n\) 次方根并吸收误差，即得
\[
r\ge (n!)^{1/n}\,2^m\,(1+o(1)).
\]
证毕。
\end{proof}

\begin{theorem}[单整数非线性全息与线性矩全息的模型内维数分离]\label{thm:spg-single-integer-vs-linear-moment-holography-gap}
固定 \(m,n\ge 1\)，并令 \(\mathcal U_{m,n}\) 为全部分辨率 \(m\) 的 dyadic polycube，\(\mathcal V_m\) 为定理 \ref{thm:spg-dyadic-finite-moment-completeness} 中的分段常数函数空间。
则：
\begin{enumerate}
  \item 边界 Gödel 全息映射
        \[
        H_m:\mathcal U_{m,n}\to \NN
        \]
        在 \(\mathcal U_{m,n}\) 上为单射；
  \item 任意有限线性观测族
        \[
        \Lambda_m=\{\ell_1,\dots,\ell_L\}\subset \mathcal V_m^\ast
        \]
        若其测量映射
        \[
        \mathcal E_{\Lambda_m}:\mathcal V_m\to\RR^L
        \]
        在整个 \(\mathcal V_m\) 上单射，则必有
        \[
        L\ge 2^{mn}.
        \]
\end{enumerate}
因此，在固定分辨率的 dyadic 全息模型内，dyadic polycube 类允许以单个非线性算术标量实现完备编码，而线性全空间重构则必然需要 \(2^{mn}\) 级坐标数。
\leanverified{paper\_spg\_single\_integer\_vs\_linear\_moment\_holography\_gap}
\leanverified{paper\_spg\_single\_integer\_vs\_linear\_moment\_gap\_seeds}
\end{theorem}
\begin{proof}
第 \(1\) 条由推论 \ref{cor:spg-boundary-godel-finite-moment-completeness} 直接给出。
第 \(2\) 条正是定理 \ref{thm:spg-linear-moment-holography-minimal-dimension} 的第 \(1\) 条。
二者并列即可得到所述模型内分离。证毕。
\end{proof}

\begin{remark}[线性维数下界的适用域]\label{rem:spg-single-integer-vs-linear-moment-holography-gap-scope}
上述线性维数下界讨论的是“在线性空间 \(\mathcal V_m\) 上完备”的观测族，而不是对有限集合 \(\mathcal U_{m,n}\) 上任意外加算术权重的编码方案作普适否定。
因此，该分离结论应理解为：在本文固定的 dyadic 全息模型中，非线性算术全息与线性全空间重构的资源门槛严格分离，而非对任意有限集合编码方式的一般性下界。
\end{remark}

\noindent
因此，固定分辨率 dyadic polycube 的线性全息并不存在低于 \(2^{mn}\) 的观测维数实现；而在各向同性总次数模型中，即便只要求达到线性完备的必要门槛，次数参数也必须直接跃迁到 \(2^m\) 量级。这表明 Prouhet--Thue--Morse 障碍并非局部低阶现象，而是与线性维数下界共同指向同一指数级资源门槛。

% --- Distillation writeback ---
% --- Distillation writeback: Wang-Zahl convex-container support audit ---
\begin{definition}[结果 17：dyadic support-separating container audit]\label{distill:wang_zahl-convex_container_non_concentration_implies_maximal_support-spg-support-separating-datum}
\textbf{日期时间：2026-04-22T17:42:33+08:00. 依赖状态：construction.}
固定定理 \ref{thm:spg-linear-moment-holography-minimal-dimension} 的分辨率参数 \(m,n\ge 1\)，令 \(N=2^m\)，并记
\[
\mathcal G_{m,n}:=\{0,\ldots,N-1\}^n
\]
为 dyadic 小立方指标网格。对 \(A\subset\mathcal G_{m,n}\)，记
\[
f_A:=\sum_{\mathbf k\in A}\ind_{C_{\mathbf k}}\in\mathcal V_m
\]
为其 active union 的示性函数。一个有限 container-budget 审计资料由四元组
\[
\mathfrak S=(\mathfrak C_m,B,S,\mathscr A_m)
\]
组成，其中 \(\mathfrak C_m\subset 2^{\mathcal G_{m,n}}\) 是有限 admissible dyadic container 族，\(B=(B_{\mathcal K})_{\mathcal K\in\mathfrak C_m}\) 是整数预算，\(1\le S\le N^n\) 是支撑门槛，\(\mathscr A_m\subset 2^{\mathcal G_{m,n}}\) 是待审计的 active supports 类。称 \(\mathfrak S\) 对 \(\mathscr A_m\) 是 \(S\)-support-separating 的，若对每个非空 \(A\in\mathscr A_m\) 满足 \(|A|<S\) 时，存在 \(\mathcal K\in\mathfrak C_m\) 使
\[
|A\cap\mathcal K|>B_{\mathcal K}.
\]
称 \(A\) 满足该资料的 Wolff-style container budget，若对全部 \(\mathcal K\in\mathfrak C_m\) 有
\[
L_{\mathcal K}(f_A)\le B_{\mathcal K},
\]
其中 \(L_{\mathcal K}\) 为定义 \ref{def:distill-spg-container-budget-functional} 的 container occupancy 泛函。此定义只记录有限 failure-extraction 审计假设；它本身不声称 admissible containers 具有内禀凸性单调、尺度继承或 rescaling closure。
\end{definition}
% --- End distillation writeback ---

% --- Distillation writeback: Wang-Zahl maximal support extraction ---
\begin{theorem}[结果 18：container budgets force finite-resolution maximal support and localize every failure]\label{distill:wang_zahl-convex_container_non_concentration_implies_maximal_support-spg-budget-forces-support}
\textbf{日期时间：2026-04-22T17:42:33+08:00. 依赖状态：rigidity.}
设 \(\mathfrak S=(\mathfrak C_m,B,S,\mathscr A_m)\) 是定义 \ref{distill:wang_zahl-convex_container_non_concentration_implies_maximal_support-spg-support-separating-datum} 意义下的 \(S\)-support-separating dyadic container audit datum。若 \(A\in\mathscr A_m\) 满足全部 Wolff-style container budgets
\[
L_{\mathcal K}(f_A)\le B_{\mathcal K}\qquad(\mathcal K\in\mathfrak C_m),
\]
则
\[
|A|\ge S.
\]
特别地，当 \(S=N^n-\Delta\) 时，active union 至多遗漏 \(\Delta\) 个 \(m\)-级 dyadic 小立方。

反之，若 \(A\in\mathscr A_m\) 非空且 \(|A|<S\)，则 overloaded-container extractor
\[
\mathfrak O(A):=\{\mathcal K\in\mathfrak C_m:\ L_{\mathcal K}(f_A)>B_{\mathcal K}\}
\]
非空；任取使 \(L_{\mathcal K}(f_A)-B_{\mathcal K}\) 最大的 \(\mathcal K_\ast\in\mathfrak O(A)\)，该 \(\mathcal K_\ast\) 即为支撑失败的具体过载 container。因此在该有限审计口径下，maximal support 的失败不可能表现为 diffuse obstruction。若完整 tensor moment box \(\mathcal M_m^{\Box}\) 可用，则 \(\mathcal K_\ast\) 的过载不等式可由命题 \ref{prop:distill-spg-container-budget-linear-audit} 离线核验。
\end{theorem}
\begin{proof}
先计算 container occupancy。由定义 \ref{def:distill-spg-container-budget-functional}，对任意 \(A\subset\mathcal G_{m,n}\) 与 \(\mathcal K\subset\mathcal G_{m,n}\)，有
\[
L_{\mathcal K}(f_A)
=2^{mn}\int_{I^n}\left(\sum_{\mathbf j\in A}\ind_{C_{\mathbf j}}(x)\right)\ind_{\mathcal K}(x)\,\dd x.
\]
不同 dyadic 小立方两两不交，且每个小立方体积为 \(2^{-mn}\)，故上式等于
\[
\#\{\mathbf j\in A:\mathbf j\in\mathcal K\}=|A\cap\mathcal K|.
\]
若 \(|A|<S\)，support-separating 条件给出某个 \(\mathcal K\in\mathfrak C_m\) 使 \(|A\cap\mathcal K|>B_{\mathcal K}\)。由刚才的等式，这正是
\[
L_{\mathcal K}(f_A)>B_{\mathcal K},
\]
与全部预算通过矛盾。因此预算全部通过时必有 \(|A|\ge S\)。同一论证也说明当 \(|A|<S\) 时集合 \(\mathfrak O(A)\) 非空；有限性保证最大 margin container \(\mathcal K_\ast\) 存在，并且它给出可逐项检查的过载不等式。最后，命题 \ref{prop:distill-spg-container-budget-linear-audit} 已说明每个 \(L_{\mathcal K}\) 都可由完整 tensor moment box 唯一线性读出，故 \(L_{\mathcal K_\ast}(f_A)>B_{\mathcal K_\ast}\) 可离线核验。该结论只使用有限 support-separating 假设与 occupancy 恒等式；并未推出 admissible containers 的内禀凸性单调或跨尺度 rescaling 继承。证毕。
\end{proof}
% --- End distillation writeback ---

% --- End distillation writeback ---

\endinput


\endinput


\begin{corollary}[单项式矩的抗扰动误差界：由 \texorpdfstring{$\gcd$}{gcd} 显式读出]\label{cor:spg-godel-moment-robust-gcd-bound}
在定理 \ref{thm:spg-godel-boundary-gcd-lipschitz-stability} 的设定下，固定 \(\alpha=(\alpha_1,\dots,\alpha_n)\in\NN^n\) 并令
\[
\omega_{\alpha}^{(1)}:=\frac{x^{\alpha+e_1}}{\alpha_1+1}\,\dd x_2\wedge\cdots\wedge \dd x_n,
\qquad
\dd\omega_{\alpha}^{(1)}=x^\alpha\,\dd x_1\cdots\dd x_n.
\]
则对任意 dyadic polycube \(U,V\subset I^n\) 有
\[
\left|\int_U x^\alpha\,\dd x_1\cdots\dd x_n-\int_V x^\alpha\,\dd x_1\cdots\dd x_n\right|
=\left|\int_{\partial U-\partial V}\omega_{\alpha}^{(1)}\right|
\le \frac{2^{-m(n-1)}}{\alpha_1+1}\,\operatorname{Dist}_m(U,V),
\]
其中 \(\operatorname{Dist}_m(U,V)\) 由定理 \ref{thm:spg-godel-boundary-gcd-lipschitz-stability} 中的 \(\gcd\)-算术表达显式给出。
\end{corollary}
\begin{proof}
由 Stokes 定理得
\(
\int_U x^\alpha\,\dd x-\int_V x^\alpha\,\dd x=\int_{\partial U-\partial V}\omega_{\alpha}^{(1)}
\)。
对任意 dyadic 取向 \((n-1)\)-面 \(\vec f\)，有 \(|x^{\alpha+e_1}|\le 1\) 且 \(\vec f\) 的 \((n-1)\)-测度为 \(2^{-m(n-1)}\)，从而
\[
\left|\int_{\vec f}\omega_{\alpha}^{(1)}\right|
\le \frac{1}{\alpha_1+1}\int_{\vec f}1\,\dd S
=\frac{2^{-m(n-1)}}{\alpha_1+1}.
\]
将 \(\partial U-\partial V\) 分解为取向面之和并用三角不等式得到
\[
\left|\int_{\partial U-\partial V}\omega_{\alpha}^{(1)}\right|
\le \frac{2^{-m(n-1)}}{\alpha_1+1}\bigl\|\partial_n[U]-\partial_n[V]\bigr\|_{\ell^1}
=\frac{2^{-m(n-1)}}{\alpha_1+1}\,\operatorname{Dist}_m(U,V),
\]
其中最后一步用定理 \ref{thm:spg-godel-boundary-gcd-lipschitz-stability}。证毕。
\end{proof}

\begin{theorem}[dyadic 边界像的 LDPC 码结构：局部校验子为共面四元关系]\label{thm:spg-dyadic-boundary-image-ldpc}
令
\[
\mathcal C_{m,n}:=\mathrm{im}\bigl(\partial_n:C_n(\mathcal Q_m;\FF_2)\to C_{n-1}(\mathcal Q_m;\FF_2)\bigr)
\]
为 dyadic 顶维 \(0/1\) 链的模 \(2\) 边界像集合。
则：
\begin{enumerate}
  \item \(\mathcal C_{m,n}\) 为长度
        \[
        N_{m,n}:=\#\mathcal F_{n-1}(\mathcal Q_m)=n(2^m+1)\,2^{m(n-1)}
        \]
        的线性码，且
        \[
        \dim \mathcal C_{m,n}=2^{mn},
        \qquad
        R_{m,n}:=\frac{\dim \mathcal C_{m,n}}{N_{m,n}}
        =\frac{2^m}{n(2^m+1)}\xrightarrow[m\to\infty]{}\frac1n.
        \]
  \item 以 \(\partial_{n-1}\) 作为校验算子，则每一条校验方程恰涉及 \(4\) 个 \((n-1)\)-胞腔坐标；因此 \(\mathcal C_{m,n}\) 构成一族规则 LDPC 码，其校验度为 \(4\)，变量度为 \(2(n-1)\)。
\end{enumerate}
\end{theorem}
\begin{proof}
\textbf{(1)}
\(\mathcal C_{m,n}\) 为 \(\FF_2\)-线性子空间由定义即得。
计数 \(N_{m,n}\) 来自：每个坐标方向有 \(2^m+1\) 个切片超平面，每个切片上有 \(2^{m(n-1)}\) 个 \((n-1)\)-胞腔，共 \(n\) 个方向。
又由定理 \ref{thm:spg-dyadic-cubical-boundary-injective}，\(\partial_n\) 在顶维链群上单射，故
\[
\dim \mathcal C_{m,n}=\rank(\partial_n)=\dim C_n(\mathcal Q_m;\FF_2)=2^{mn}.
\]
码率公式随即得到。
\par
\textbf{(2)}
在立方格点局部，一个 \((n-2)\)-胞腔位于两条坐标方向的交叉处，恰被四个 \((n-1)\)-胞腔围绕；因此 \(\partial_{n-1}\) 的对应一行支撑大小为 \(4\)。
另一方面，每个 \((n-1)\)-胞腔的边界含 \(2(n-1)\) 个 \((n-2)\)-胞腔，故变量度为 \(2(n-1)\)。证毕。
\end{proof}

\begin{corollary}[dyadic 边界像的最小距离]\label{cor:spg-dyadic-boundary-image-min-distance}
在定理 \ref{thm:spg-dyadic-boundary-image-ldpc} 的记号下，以 Hamming 重量
\[
\mathrm{wt}(b):=\#\{\sigma\in\mathcal F_{n-1}(\mathcal Q_m):\ b(\sigma)=1\}
\qquad\bigl(b\in C_{n-1}(\mathcal Q_m;\FF_2)\bigr)
\]
定义 \(\mathcal C_{m,n}\) 的最小距离
\(
d_{\min}(\mathcal C_{m,n}):=\min\{\mathrm{wt}(b):b\in\mathcal C_{m,n}\setminus\{0\}\}
\)。
则对一切 \(m\ge 1\) 与 \(n\ge 1\) 有
\[
\boxed{\ d_{\min}(\mathcal C_{m,n})=2n\ }.
\]
\end{corollary}
\begin{proof}
取任一 \(n\)-胞腔 \(Q\in\mathcal F_n(\mathcal Q_m)\)，则 \(\partial_n Q\in\mathcal C_{m,n}\setminus\{0\}\) 且其支撑恰由 \(Q\) 的 \(2n\) 张 \((n-1)\)-面构成，故 \(d_{\min}(\mathcal C_{m,n})\le 2n\)。
\par
反之，取任意非零 \(b=\partial_n x\in\mathcal C_{m,n}\) 且 \(x\in C_n(\mathcal Q_m;\FF_2)\) 非零。
把 \(x\) 的支撑视为非空 \(n\)-胞腔集合 \(S\subset\mathcal F_n(\mathcal Q_m)\)。
对每个坐标方向 \(i\in\{1,\dots,n\}\)，取 \(S\) 中在第 \(i\) 坐标上达到最小值的一个胞腔 \(Q_i^{-}\)；
则其 \((-e_i)\)-外侧面不可能与另一胞腔相邻（否则第 \(i\) 坐标还能更小），因而在 \(\partial_n x\) 中出现。
同理取达到最大值的一个胞腔 \(Q_i^{+}\) 得到一张 \((+e_i)\)-外侧面亦在 \(\partial_n x\) 中出现。
对不同的 \(i\) 与不同符号得到的外侧面两两不同，故 \(\mathrm{wt}(b)\ge 2n\)。
于是 \(d_{\min}(\mathcal C_{m,n})\ge 2n\)，合并得等号。证毕。
\end{proof}

\endinput

